
\section{射影概形的积：\trSegre{}嵌入}

我们接下来描述射影 $A$-概形在 $A$ 上的积。（最初最令人感兴趣的情形是 $A=k$。）为此，只需描述 $\mathbb{P}^m_A\times_A\mathbb{P}^n_A$，因为任何\trProjectiveScheme{}都\trClosedEmbedding{}于某个 $\mathbb{P}^m_A$，而\trClosedEmbedding{}在\trBaseChange{}下表现良好，于是若 $X\hookrightarrow\mathbb{P}^m_A$ 与 $Y\hookrightarrow\mathbb{P}^n_A$ 都是\trClosedEmbedding，则 $X\times_A Y\hookrightarrow\mathbb{P}^m_A\times_A\mathbb{P}^n_A$ 也是\trClosedEmbedding，由 $X$ 与 $Y$ 的方程割出（\S10.2(3)）。我们将描述 $\mathbb{P}^m_A\times_A\mathbb{P}^n_A$，并看到它同样是一个射影 $A$-\trScheme。（于是，若 $X$ 与 $Y$ 是射影 $A$-\trScheme，则它们在 $A$ 上的积 $X\times_A Y$ 也是射影 $A$-\trScheme。）

在做这件事之前，我们先在一个特殊情形里，从经典射影空间（非零向量模去非零标量，\cref{ex:closed-points-pn-traditional}）获得一点动机。我们的\trMap{}将把 $[x_0,x_1,x_2]\times[y_0,y_1]$ 送到 $\mathbb{P}^5$ 中的一个点，我们把它的诸\trCoordinate{}设想成下面这张「乘法表」
\[
\left[\begin{array}{ccc}
x_0y_0, & x_1y_0, & x_2y_0, \\
x_0y_1, & x_1y_1, & x_2y_1
\end{array}\right]
\]
中的表项。

这确实是一个\trWellDefined{}的集合\trMap。注意：所得矩阵的\trRank{}为 $1$，而且从这个矩阵出发，我们可以在相差非零标量倍数的意义下读出 $[x_0,x_1,x_2]$ 与 $[y_0,y_1]$。例如，要读出点 $[x_0,x_1,x_2]\in\mathbb{P}^2$，就取第一行，除非它全为零，此时便取第二行。（两行不可能都全为零。）总而言之：在经典射影几何中，给定了 $\mathbb{P}^m$ 与 $\mathbb{P}^n$ 的点，我们就造出了 $\mathbb{P}^{mn+m+n}$ 中的一个点；而从这个 $\mathbb{P}^{mn+m+n}$ 中的点出发，我们又能恢复出 $\mathbb{P}^m$ 与 $\mathbb{P}^n$ 的点。

动机已足，我们回到代数几何。定义一个\trMap
\[
\mathbb{P}^m_A\times_A\mathbb{P}^n_A\longrightarrow\mathbb{P}^{mn+m+n}_A
\]
如下：
\[
\begin{aligned}
\bigl([x_0,\dots,x_m],[y_0,\dots,y_n]\bigr)&\longmapsto[z_{00},z_{01},\dots,z_{ij},\dots,z_{mn}]\\
&=[x_0y_0,\ x_0y_1,\ \dots,\ x_iy_j,\ \dots,\ x_my_n].
\end{aligned}
\]
说得更显式些：考虑从仿射开集 $U_i\times V_j$（其中 $U_i=D(x_i)$、$V_j=D(y_j)$）到仿射开集 $W_{ij}=D(z_{ij})$ 的映射
\[
(x_{0/i},\dots,x_{m/i},y_{0/j},\dots,y_{n/j})\longmapsto(x_{0/i}y_{0/j},\dots,x_{i/i}y_{j/j},\dots,x_{m/i}y_{n/j}),
\]
或者用代数的语言说，$z_{ab/ij}\mapsto x_{a/i}y_{b/j}$。

\textbf{10.6.A.（练习）}\quad 检验：这些映射粘合起来，给出一个\trWellDefined{}的\trMorphism{} $\mathbb{P}^m_A\times_A\mathbb{P}^n_A\to\mathbb{P}^{mn+m+n}_A$。

\textbf{10.6.1.}\quad 我们接下来证明：这个\trMorphism{}是一个\trClosedEmbedding。这可以在目标的一个\trOpenCover{}上检验（「是\trClosedEmbedding」这一观念是\trAffineLocal{}的，\mbox{练习~9.1.D}）。让我们在 $z_{ij}\neq 0$ 的那个开集上检验。这个开集在 $\mathbb{P}^m_A\times_A\mathbb{P}^n_A$ 中的原像，是 $x_i\neq 0$ 且 $y_j\neq 0$ 的\trLocus，也就是 $U_i\times V_j$。如上所述，\trRingMorphism{}由 $z_{ab/ij}\mapsto x_{a/i}y_{b/j}$ 给出；这显然是\trSurjectiveMap{}的，因为 $z_{aj/ij}\mapsto x_{a/i}$、$z_{ib/ij}\mapsto y_{b/j}$。（这一特设描述的推广将在\mbox{练习~15.2.E}中给出。）

这个\trMap{}称为\textbf{\trSegre{}态射}（Segre morphism）或\textbf{\trSegre{}嵌入}（Segre embedding）。若 $A$ 是一个\trField，其像便称为\textbf{\trSegre{}簇}（Segre variety）。

\textbf{10.6.B.（练习）}\quad 证明：\trSegre{}概形（\trSegre{}嵌入的像）由对应于下式的方程\trSchemeTheoretic{}地割出：
\[
\operatorname{rank}\begin{pmatrix}
z_{00} & \cdots & z_{0n}\\
\vdots & \ddots & \vdots\\
z_{m0} & \cdots & z_{mn}
\end{pmatrix}=1,
\]
也就是说，一切 $2\times 2$ \trMinor{}都消失。提示：设你有一个关于诸 $z_{ij}$ 的多项式，它在代入 $z_{ij}=x_iy_j$ 后变为零。给出一个配方：逐个减去形如「\trMonomial{}乘 $2\times 2$ \trMinor」的多项式，使得最终结果为 $0$。（在特殊情形下，\trVeroneseEmbedding{}的类似问题正是\mbox{练习~9.3.I}与\mbox{练习~9.3.K}的内容。）

\textbf{10.6.2.\ 重要例子。}\quad 让我们考虑第一个非平凡的例子：$m=n=1$。我们得到 $\mathbb{P}^1_k\times_k\mathbb{P}^1_k\hookrightarrow\mathbb{P}^3_k$（$k$ 是一个\trField）。我们得到单独一个方程
\[
\operatorname{rank}\begin{pmatrix} z_{00} & z_{01}\\ z_{10} & z_{11}\end{pmatrix}=1,
\]
即 $z_{00}z_{11}-z_{01}z_{10}=0$。我们又见到了老朋友——\trQuadricSurface{}（\S9.3.9）！于是：光滑\trQuadricSurface{} $wz-xy=0$（图~\ref{fig:ruling-quadric-surface}）同构于 $\mathbb{P}^1_k\times_k\mathbb{P}^1_k$。回想\mbox{练习~9.3.L}：这个二次曲面有两族（\trRuling）直线。你或许愿意检验一下：其中一族直线对应于 $\{x\}\times_k\mathbb{P}^1_k$ 的像随 $x$ 的变动，另一族对应于 $\mathbb{P}^1_k\times_k\{y\}$ 的像随 $y$ 的变动。

若 $k$ 是\trCharacteristic{}非 $2$ 的\trAlgClosedField，则由二次型的可\trDiagonalize{}性（\mbox{练习~5.4.J}），四个变量上全体\trRank{}为 $4$（「\trFullRank」）的\trQuadraticForm{}彼此同构，于是\trCharacteristic{}非 $2$ 的\trAlgClosedField{}上全体\trRank{}为 $4$ 的二次曲面都同构于 $\mathbb{P}^1_k\times_k\mathbb{P}^1_k$。

注意：在一个非\trAlgClosed{}的域上，这不再成立。例如在 $\mathbb{R}$ 上，$w^2+x^2+y^2+z^2=0$（在 $\mathbb{P}^3_\mathbb{R}$ 中）并不同构于 $\mathbb{P}^1_\mathbb{R}\times_\mathbb{R}\mathbb{P}^1_\mathbb{R}$。理由：前者没有实点（real point），而后者的实点多得很。

接下来的两道练习，你可以按任意顺序做。第二道可以用来证明第一道，但第一道或许能让你对第二道有所领悟。

\textbf{10.6.C.（练习：\trSegre{}嵌入的无坐标（coordinate-free）描述）}\quad 证明：\trSegre{}嵌入可以解释为 $\mathbb{P}V\times\mathbb{P}W\to\mathbb{P}(V\otimes W)$，经由分次环的满射
\[
\Sym^\bullet(V^\vee\otimes W^\vee)\twoheadrightarrow\bigoplus_{t=0}^{\infty}\bigl(\Sym^t V^\vee\bigr)\otimes\bigl(\Sym^t W^\vee\bigr)
\]
「朝相反的方向」给出。

\textbf{10.6.D.（练习：射影 $A$-概形之积的一般无坐标描述）}\quad 设 $\mathsf{S}_\bullet$ 与 $\mathsf{T}_\bullet$ 是 $A$ 上\trFinitelyGenerated{}的\trGradedRing。描述一个同构
\[
(\Proj \mathsf{S}_\bullet)\times_A(\Proj \mathsf{T}_\bullet)\overset{\sim}{\longleftrightarrow}\Proj\bigoplus_{n=0}^{\infty}\bigl(\mathsf{S}_n\otimes_A\mathsf{T}_n\bigr)
\]
（但愿分次环 $\bigoplus_{n=0}^{\infty}(\mathsf{S}_n\otimes_A\mathsf{T}_n)$ 中乘法的定义是清楚的）。
